% !TEX root = ../main.tex


\section{Order-Flow Data and Tie Frequency}
\label{app:zq-data}

\subsection{Nasdaq Sampling and Replay}
\label{app:zq-data-replay}

\paragraph{Day selection.}
We reconstruct the AAPL book from Nasdaq TotalView-ITCH order and trade messages. The sampling procedure was defined before measuring tie frequency; an extended event screen was applied after the first replay without changing the four analysis dates. Of 219 trading days from 2 January through 13 November 2026, 119 remained after excluding days around FOMC decisions, CPI and employment releases, option expiries, index reconstitution, the election, early closes, and holidays. We shuffled these days with a fixed seed and then applied a further screen for macro releases, central-bank minutes, month ends, AAPL and NVDA earnings and the following day, and AAPL ex-dividend and product events. We also excluded days with an AAPL return beyond two standard deviations of the preceding 60 days, volume outside one-half to twice the 20-day median, a range above twice its median, or shocks to six major technology stocks and selected semiconductor, oil, Treasury, or gold funds. Sixty-two days passed this second screen. We used the first four in shuffled order: 23 January (3,176,565 order messages), 4 May (5,629,415), 9 July (6,065,233), and 20 July 2026 (11,064,805). We separately replayed 27 August (5,092,753 messages), rejected because it followed NVDA earnings; it is \emph{not} included in frequency estimates.

\paragraph{Book reconstruction.}
We partition the regular session (09:30--16:00 New York time) into non-overlapping 0.1-, 1-, 10-, or 60-second batches on a one-cent price grid. At each tick, bid volume comprises resting bids at the end of the batch, displayed bids executed during the batch, and buyer-initiated trades; ask volume is constructed symmetrically. We exclude cancelled interest, round sub-cent bid prices down and ask prices up, and omit crosses and non-displayed executions whose aggressor side is unavailable (24--33\% of trade messages). Each price window extends 300 ticks beyond the batch's trade prices on both sides. Resting depth is essential: executions alone yield a plateau in nearly every batch, whereas most books reconstructed with depth have a unique maximizing tick.

\paragraph{Checks and benchmark selection.}
Reconstructed resting depth was never negative. Across 156,950 crossing batches on the four analysis days, no maximizing plateau reached a window boundary and each batch had a candidate clearing tick. The Python replay also passed test cases from the Rust implementation. For each of the five replayed days and four durations, we selected the first plateau with one minimum-imbalance tick and the first with several, where available. This produced 32 benchmark books with plateau widths of 2--35 ticks; twelve also had a tie after imbalance minimization. Only the four screened days enter the frequency denominators.

\subsection{Tie-Frequency Variation}
\label{app:zq-data-ties}

The main text reports the mean tie frequencies at all four batch lengths. The one-second plateau rate averaged 3.63\% across the four screened days, with a 95\% Student-$t$ interval of 1.5--5.7\% computed from \emph{four daily observations}, not from individual batches. The daily rates below distinguish any volume-maximizing plateau from a tie remaining after imbalance minimization.

\begin{center}
\small
\begin{tabular}{@{}lrr@{}}
\toprule
Day (2026) & Plateau (\%) & Several imbalance minima (\%) \\
\midrule
23 January & 2.16 & 0.07 \\
4 May & 3.62 & 0.13 \\
9 July & 3.38 & 0.35 \\
20 July & 5.37 & 0.75 \\
27 August (rejected) & 3.33 & 0.48 \\
\bottomrule
\end{tabular}
\end{center}

Across all four batch lengths, 546 crossing batches had several ticks tied on minimum imbalance. In 462 of those, two ticks also satisfied the at-price volume condition; the rule selected the lower one. At durations up to ten seconds, plateaus were more frequent in the first half hour of trading than in the last. The rejected day's plateau rates fell within the range observed on the screened days at every duration. %Muni Toke studies the clearing-price range under a stochastic call-auction model \cite{munitoke2015}; our frequencies instead describe clearing reconstructed continuous-market order flow under our rule.

\subsection{Synthetic-Book Calibration}
\label{app:zq-data-calibration}

We centre each crossing book on its clearing tick and estimate tick occupancy and volume quantiles within 100 ticks of the centre and in twelve wider bands extending to 10,000 ticks. We retain joint volumes within ten ticks and total bid depth above and ask depth below that region. We tested three generators on 20,000 books of 601 ticks per batch duration: \emph{independent ticks} sampled occupancy and size separately at each price; \emph{core occupancy} copied the occupancy pattern within five ticks from a replayed batch but sampled sizes independently; and the \emph{bootstrap} copied joint volumes within ten ticks, sampled the periphery from the marginals, and rescaled peripheral totals to match the source batch.

\noindent\textit{One-second batches, percentage of crossing books (20,000 books per synthetic generator):}
\begin{center}
\small
\setlength{\tabcolsep}{4pt}
\begin{tabular}{@{}lrrrr@{}}
\toprule
Source & Unique max. & Plateau, one min. & Plateau, several & Median $V_{max}$ \\
\midrule
Replay (pooled) & 96.21 & 3.43 & 0.36 & 120 \\
Independent ticks & 85.40 & 12.45 & 2.15 & 132 \\
Core occupancy & 93.89 & 5.91 & 0.20 & 119 \\
Bootstrap & 95.94 & 3.74 & 0.32 & 120 \\
\bottomrule
\end{tabular}
\end{center}

The bootstrap's total one-second plateau frequency is 4.06\%, compared with 3.79\% in the pooled replay; its frequencies also closely match the other three batch durations. Independent tick sampling substantially overstates ties. We therefore draw the 25 calibrated benchmark books from the one-second bootstrap, five each at 100, 250, 500, 1,000, and 2,000 ticks. This calibration tests reproduction of the measured distribution, not how traders would behave in a deployed frequent batch auction.

\subsection{Periodic-Auction Data from Cboe Europe}
\label{app:zq-cboe}

We additionally use 111,871 BXE messages for Orange SA (ORAp) on 1 September 2020. The feed contains 922 periodic-auction updates (493 with a positive indicative price), 362 auction summaries with final price and volume, and 400 executions flagged as auction trades. We associate each positive-price update with the next summary and executions with a summary at the same timestamp. The execution quantities sum to the reported size in all 362 auctions. The 362 auctions traded 238,921 shares, with a median of 433.5 shares per auction; 327 had one execution, 32 had two, and three had three. For 320 auctions, at most 100 ms elapsed from the first assigned positive-price update to the summary; for 35, more than one second elapsed. These observed intervals need not equal auction durations.

The feed does not reveal individual auction orders, including unexecuted interest. To test proofs against observed prices and volumes, we therefore \emph{construct} a 101-tick book on a EUR 0.001 grid centred at each reported price, put the executed volume on both sides at the centre, and set all other ticks to zero. Both implementations proved all 362 constructed books and matched each published price and volume. These tests do not verify the venue's actual private book or its tie-breaking: our constructed books have a unique maximizing tick by design. The venue's auctions also occur at irregular times, unlike fixed-interval batches.

\paragraph{Limits.}
The Nasdaq sample covers one stock on four analysis days and includes no IO orders. The Cboe sample covers one stock on one day; its message interpretation requires checking against the specification applicable in 2020. Publishing replayed or calibrated books may be restricted by the feeds' data licences.
